\documentclass[12pt,a4paper]{ctexart}
\usepackage[margin=1in]{geometry}
\usepackage{booktabs}
\usepackage{longtable}
\usepackage{array}
\usepackage{tabularx}
\usepackage{hyperref}

\title{研一下第一周组会汇报}
\author{ }
\date{\today}

\begin{document}
\maketitle

\section{本周基线定义}
本周将 \texttt{src/train\_balanced\_sampling\_sep\_mlp.py} 的 \textbf{concat + NET readout} 版本作为基线。
其中 HGAT 关键配置口径为：
\begin{itemize}
  \item 多层 HGAT（\texttt{hgat\_layers}）与层内 dropout（\texttt{hgat\_dropout}）；
  \item readout 支持 NET 统计特征（\texttt{hgat\_use\_net\_readout}）；
  \item 可选节点类型注意力融合（\texttt{hgat\_type\_attn\_readout}，当前实验中可关闭）；
  \item arc 条件化主线采用 \texttt{concat} 模式（即将 arc embedding 与 \(x,z\) 拼接回归）。
\end{itemize}

\section{HGAT 本周改动点（相对早期版本）}
\subsection{编码器结构}
\begin{itemize}
  \item 引入块级残差 + LayerNorm + Dropout（\texttt{HGATBlock}）。
  \item 编码器支持可配置层数（\texttt{num\_layers}）和 dropout。
\end{itemize}

\subsection{图特征与读出}
\begin{itemize}
  \item NET 特征从简单占位增强为：\texttt{is\_vdd/is\_vss/is\_top\_pin/deg\_norm}。
  \item MOS 特征保持 \(\log W\) 与 \(\log L\)。
  \item 读出由单均值扩展为统计读出（mean/max/std），并支持是否纳入 NET 读出。
\end{itemize}

\subsection{训练脚本接线统一}
README 中涉及的训练脚本均已统一为同一 HGAT 初始化口径（支持上述参数）：
\begin{itemize}
  \item \texttt{src/train.py}
  \item \texttt{src/train\_base\_pt\_ft.py}
  \item \texttt{src/train\_balanced\_sampling.py}
  \item \texttt{src/train\_balanced\_sampling\_sep\_mlp.py}
  \item \texttt{src/train\_balanced\_sampling\_sep\_mlp\_disentangle\_cell\_type.py}
  \item \texttt{src/train\_balanced\_sampling\_sep\_mlp\_disentangle\_domain.py}
  \item \texttt{src/train\_balanced\_sampling\_sep\_mlp\_bayesian.py}
\end{itemize}

\begin{itemize}
  \item \texttt{src/train\_balanced\_sampling\_sep\_mlp\_shared\_calib.py}
\end{itemize}
对齐到同一 HGAT 版本。

\section{脚本差异说明}
\subsection{\texttt{shared\_calib} 相对 \texttt{sep\_mlp}}
\texttt{src/train\_balanced\_sampling\_sep\_mlp\_shared\_calib.py} 相对
\texttt{src/train\_balanced\_sampling\_sep\_mlp.py} 的核心差异：
\begin{itemize}
  \item 回归头结构改为共享主干 + 域校正（calibration）分支，而非完全分离的双头 MLP。
  \item 训练目标仍是 src/tgt 平衡采样回归，但通过共享与校正分工来减少参数冗余。
  \item HGAT 编码器、图缓存与 freeze/unfreeze 流程与主线脚本保持一致口径。
\end{itemize}

\subsection{\texttt{disentangle\_domain} 相对 \texttt{sep\_mlp}}
\texttt{src/train\_balanced\_sampling\_sep\_mlp\_disentangle\_domain.py} 在主线回归损失外，增加解耦约束：
\begin{itemize}
  \item 将表示拆分为 design/process 子空间；
  \item 对 process 子空间施加域相关约束（如对比学习项），对 design 子空间施加分布对齐项（如 CMD）；
  \item 总损失 = 回归损失 + \(\lambda_{clr}\) 对比项 + \(\lambda_{cmd}\) 对齐项。
\end{itemize}

\subsection{\texttt{disentangle\_domain} 相对 \texttt{sep\_mlp}}
\texttt{src/train\_balanced\_sampling\_sep\_mlp\_disentangle\_domain.py} 相对
\texttt{src/train\_balanced\_sampling\_sep\_mlp.py} 的核心修改为在主线回归上叠加域解耦约束：
\begin{itemize}
  \item 模型输出从单一回归分支扩展为回归预测 + 表示分解（design/process）分支；
  \item 在每个 batch 中基于 src/tgt 配对，增加域相关对比约束（CLR）与分布对齐约束（CMD）；
  \item 总损失从纯回归损失扩展为
  \(\mathcal{L}_{reg} + \lambda_{clr}\mathcal{L}_{clr} + \lambda_{cmd}\mathcal{L}_{cmd}\)；
  \item 其余流程（HGAT 编码、图缓存、平衡采样主循环）与 \texttt{sep\_mlp} 主线保持一致。
\end{itemize}

\subsection{\texttt{bayesian} 相对 \texttt{sep\_mlp}}
\texttt{src/train\_balanced\_sampling\_sep\_mlp\_bayesian.py} 相对
\texttt{src/train\_balanced\_sampling\_sep\_mlp.py} 的主要修改如下：
\begin{itemize}
	\item 在回归头末层由确定性线性层改为贝叶斯线性层（参数为分布，训练时采样权重）；
	\item 在原有回归损失基础上可选加入 KL 正则（\texttt{--kl}、\texttt{--weight\_kl}、\texttt{--kl\_start\_ep}）；
	\item 验证阶段支持 MC 多次采样并取均值预测（\texttt{--mc\_num}、\texttt{--test\_bayesian}）；
	\item HGAT 编码、arc 条件化、平衡采样和数据接口与 \texttt{sep\_mlp} 主线保持一致，确保可直接同口径对照。
\end{itemize}

\section{运行结果对照表}
\renewcommand{\arraystretch}{1.2}
\begin{table}[!htbp]
	\centering
	\small
	\setlength{\tabcolsep}{10pt}
	\begin{tabular}{>{\raggedright\arraybackslash}p{0.48\textwidth}cc}
		\toprule
		模型别名 & 不冻结 $R^2$ & 冻结$R^2$ \\
		\midrule
		sampling & \underline{0.855} & \underline{0.858} \\
		sep\_mlp & \underline{0.853} & \underline{0.831} \\
		disentangle\_domain.py & \underline{0.824} & \underline{0.821} \\
		
		\bottomrule
	\end{tabular}
	\caption{本周改动HGAT之前冻结与不冻结HGAT的$R^2$对比（hgat\_layers=2）}
\end{table}

\noindent\textbf{别名与脚本映射：}
\begin{itemize}
	\item sampling $\rightarrow$ \texttt{train\_balanced\_sampling.py}
	\item sep\_mlp $\rightarrow$ \texttt{train\_balanced\_sampling\_sep\_mlp.py}
	\item disentangle\_domain $\rightarrow$ \texttt{train\_balanced\_sampling\_sep\_mlp\_disentangle\_domain.py}
	
	
\end{itemize}

\begin{table}[htbp]
\centering
\small
\setlength{\tabcolsep}{10pt}
\begin{tabular}{>{\raggedright\arraybackslash}p{0.48\textwidth}cc}
\toprule
模型别名 & 验证集 $R^2$ & 验证集 Loss \\
\midrule
sep\_mlp（concat + NET readout） & \underline{0.8621} & \underline{0.140017} \\
shared\_calib & \underline{0.8739} & \underline{0.128001} \\
disentangle\_domain & \underline{0.7774} & \underline{0.225947} \\
disentangle\_cell\_type & \underline{0.7774} & \underline{0.225947} \\
bayesian  & \underline{0.8662} & \underline{0.135781} \\           
\bottomrule
\end{tabular}
\caption{各脚本验证集结果对照（hgat\_layers=2）}
\end{table}

\noindent\textbf{别名与脚本映射：}
\begin{itemize}
  \item sep\_mlp $\rightarrow$ \texttt{train\_balanced\_sampling\_sep\_mlp.py}
  \item shared\_calib $\rightarrow$ \texttt{train\_balanced\_sampling\_sep\_mlp\_shared\_calib.py}
  \item disentangle\_domain $\rightarrow$ \texttt{train\_balanced\_sampling\_sep\_mlp\_disentangle\_domain.py}
  \item disentangle\_cell\_type $\rightarrow$ \texttt{train\_balanced\_sampling\_sep\_mlp\_disentangle\_cell\_type.py}
  \item bayesian $\rightarrow$ \texttt{train\_balanced\_sampling\_sep\_mlp\_bayesian.py}
\end{itemize}

\begin{table}[htbp]
	\centering
	\small
	\setlength{\tabcolsep}{10pt}
	\begin{tabular}{>{\raggedright\arraybackslash}p{0.48\textwidth}cc}
		\toprule
		hgat\_layers & 验证集 $R^2$ & 验证集 Loss \\
		\midrule
		1 & \underline{0.8414} & \underline{0.161019} \\
		2 & \underline{0.8621} & \underline{0.140017} \\
		3 & \underline{0.8684} & \underline{0.133551} \\
		4 & \underline{0.8690} & \underline{0.133027} \\
		5 & \underline{0.8694} & \underline{0.132526} \\
		\bottomrule
	\end{tabular}
	\caption{HGAT消息传递层数与训练结果关系（以train\_balanced\_sampling\_sep\_mlp.py为例）}
\end{table}




\end{document}
